\documentclass[11pt]{scrartcl}
\usepackage[sexy]{evan}
\usepackage{circuitikz}
%% Many thanks to Evan Chen (web.evanchen.cc) for designing evan.sty, in which all the documentation of this project was crafted.
 
\title{Left Column Top Row (LCTR)}
\subtitle{Integer Partition Combinatorial Game}
\author{Soumitro Shovon Dwip \\ Aayan Pratim Deb}

\begin{document}
	
	\maketitle
	
	\begin{mdframed}[style=mdpurplebox,frametitle={Game statement}]
		Given an integer partition of $n$, we consider the impartial combinatorial game LCTR in which moves consist of removing either the left column or top row of its Young diagram.\footnote[1]{arXiv:2207.05599 [math.CO]}. The program was written in JavaScript with aid of memoised Sprague-Grundy values. 	
	\end{mdframed} 
	$9 = 5 + 2 + 2 + 1 $ ; thus, the partition $( 5, 2^2, 1 ) \to 9$ would look like the following as the Young Diagram: 
	
	\begin{figure}[!ht]
		\centering
		\resizebox{0.3\textwidth}{!}{%
			\begin{circuitikz}
				\tikzstyle{every node}=[font=\LARGE]
				\draw [short] (1.25,12) -- (1.25,7);
				\draw [short] (1.25,7) -- (2.5,7);
				\draw [short] (2.5,7) -- (2.5,12);
				\draw [short] (1.25,12) -- (7.5,12);
				\draw [short] (7.5,12) -- (7.5,10.75);
				\draw [short] (1.25,10.75) -- (7.5,10.75);
				\draw [short] (1.25,8.25) -- (3.75,8.25);
				\draw [short] (3.75,12) -- (3.75,8.25);
				\draw [short] (1.25,9.5) -- (3.75,9.5);
				\draw [short] (5,12) -- (5,10.75);
				\draw [short] (6.25,12) -- (6.25,10.75);
			\end{circuitikz}
		}%
		\caption{Young Diagram of $( 5, 2^2, 1 )$}
		\end{figure}
		\section*{Moves} 
		There are two types of move, which both the players can perform (which makes the game an impartial one). 
		
		 \newpage 
		 \begin{itemize}
		 	
		 	\item[\textbf{(a)}] Left Column Removal
		 	
		 	\begin{figure}[!ht]
		 		\centering
		 		\resizebox{0.5\textwidth}{!}{%
		 			\begin{circuitikz}
		 				\tikzstyle{every node}=[font=\LARGE]
		 				
		 				\draw [short] (2.5,10.75) -- (2.5,4.75);
		 				\draw [short] (2.5,4.5) -- (2.5,5.25);
		 				\draw [short] (1.25,4.5) -- (2.5,4.5);
		 				\draw [short] (1.25,8.25) -- (5,8.25);
		 				\draw [short] (5,10.75) -- (5,8.25);
		 				\draw [short] (3.75,10.75) -- (3.75,5.75);
		 				\draw [short] (1.25,5.75) -- (3.75,5.75);
		 				\draw [short] (1.25,7) -- (3.75,7);
		 				\draw [short] (6.25,10.75) -- (6.25,9.5);
		 				\draw [->, >=Stealth] (8.75,8.25) -- (12.25,8.25);
		 				\draw [short] (13.75,9.5) -- (13.75,4.5);
		 				\draw [short] (13.75,4.5) -- (15,4.5);
		 				\draw [short] (13.75,9.5) -- (17.5,9.5);
		 				\draw [short] (17.5,9.5) -- (17.5,8.25);
		 				\draw [short] (13.75,8.25) -- (17.5,8.25);
		 				\draw [short] (13.75,5.75) -- (16.25,5.75);
		 				\draw [ dashed] (0,10.75) rectangle  (1.25,4.5);
		 				\draw [short] (0,9.5) -- (6.25,9.5);
		 				\draw [short] (1.25,10.75) -- (6.25,10.75);
		 				\draw [short] (0,8.25) -- (1.25,8.25);
		 				\draw [short] (0,7) -- (1.25,7);
		 				\draw [short] (0,5.75) -- (1.25,5.75);
		 				\draw [short] (13.75,9.5) -- (13.75,10.75);
		 				\draw [short] (13.75,10.75) -- (18.75,10.75);
		 				\draw [short] (18.75,10.75) -- (18.75,9.5);
		 				\draw [short] (17.5,9.5) -- (18.75,9.5);
		 				\draw [short] (15,10.75) -- (15,4.5);
		 				\draw [short] (16.25,10.75) -- (16.25,5.75);
		 				\draw [short] (17.5,10.75) -- (17.5,9.5);
		 				\draw [short] (13.75,7) -- (16.25,7);
		 			\end{circuitikz}
		 		}%
		 		
		 		\label{fig:my_label}
		 	\end{figure}
		 	\item[\textbf{(b)}] Top Row Removal
		 	 
		 	 \begin{figure}[!ht]
		 	 	\centering
		 	 	\resizebox{0.5\textwidth}{!}{%
		 	 		\begin{circuitikz}
		 	 			\tikzstyle{every node}=[font=\LARGE]
		 	 			\draw [ dashed] (1.25,10.75) rectangle  (7.5,9.5);
		 	 			\draw [short] (2.5,10.75) -- (2.5,4.75);
		 	 			\draw [short] (2.5,4.5) -- (2.5,5.25);
		 	 			\draw [short] (1.25,9.5) -- (1.25,4.5);
		 	 			\draw [short] (1.25,4.5) -- (2.5,4.5);
		 	 			\draw [short] (1.25,8.25) -- (5,8.25);
		 	 			\draw [short] (5,10.75) -- (5,8.25);
		 	 			\draw [short] (3.75,10.75) -- (3.75,5.75);
		 	 			\draw [short] (1.25,5.75) -- (3.75,5.75);
		 	 			\draw [short] (1.25,7) -- (3.75,7);
		 	 			\draw [short] (6.25,10.75) -- (6.25,9.5);
		 	 			\draw [->, >=Stealth] (8.75,8.25) -- (12.25,8.25);
		 	 			\draw [short] (13.75,9.5) -- (13.75,4.5);
		 	 			\draw [short] (13.75,4.5) -- (15,4.5);
		 	 			\draw [short] (13.75,9.5) -- (17.5,9.5);
		 	 			\draw [short] (17.5,9.5) -- (17.5,8.25);
		 	 			\draw [short] (13.75,8.25) -- (17.5,8.25);
		 	 			\draw [short] (13.75,5.75) -- (16.25,5.75);
		 	 			\draw [short] (16.25,9.5) -- (16.25,5.75);
		 	 			\draw [short] (15,9.5) -- (15,5.75);
		 	 			\draw [short] (13.75,7) -- (16.25,7);
		 	 			\draw [short] (15,5.75) -- (15,4.5);
		 	 		\end{circuitikz}
		 	 	}%
		 	 	\label{fig:my_label}
		 	 \end{figure} 
		 \end{itemize}
		 

		 
		\begin{thebibliography}{}  
			
			\bibitem[Gottlieb et al.(2024)]{gottlieb2024sprague}  
			Eric Gottlieb, Matja{\v{z}} Krnc, and Peter Mur{\v{s}}i{\v{c}}.  
			\newblock Sprague--Grundy values and complexity for {LCTR}.  
			\newblock {\em Discrete Applied Mathematics}, 346:154--169, 2024.  
			\newblock Elsevier.  
			
			\bibitem[Gottlieb et al.(2023)]{gottlieb2023some}  
			Eric Gottlieb, Jelena Ili{\'c}, and Matja{\v{z}} Krnc.  
			\newblock Some results on {LCTR}, an impartial game on partitions.  
			\newblock {\em Involve, a Journal of Mathematics}, 16(3):529--546, 2023.  
			\newblock Mathematical Sciences Publishers.  
			
		\end{thebibliography}  
		

		

	



	
	
	
	
\end{document}